We want to show that the functional

[
J(u)=\int_U |\nabla u|^2,dx-\int_U fu,dx
]

attains its minimum on (H_0^1(U)), i.e. that there exists (u_0\in H_0^1(U)) such that

[
J(u_0)=\inf_{u\in H_0^1(U)}J(u).
]

This is the standard **direct method in the calculus of variations**, one of the foundational techniques in Calculus of Variations.

---

# Step 1: Define the infimum

Let

[
m=\inf_{u\in H_0^1(U)}J(u).
]

Choose a minimizing sequence ((u_j)\subset H_0^1(U)) such that

[
J(u_j)\to m.
]

So

[
\int_U |\nabla u_j|^2,dx-\int_U fu_j,dx\to m.
]

We want to show ((u_j)) has a weakly convergent subsequence.

---

# Step 2: Show the minimizing sequence is bounded

We estimate the linear term using Cauchy–Schwarz:

[
\left|\int_U fu_j,dx\right|
\le
|f|*{L^2(U)}|u_j|*{L^2(U)}.
]

Using Poincaré:

[
|u_j|*{L^2(U)}
\le
C|\nabla u_j|*{L^2(U)}.
]

Hence

[
\left|\int_U fu_j,dx\right|
\le
C|f|*{L^2(U)}|\nabla u_j|*{L^2(U)}.
]

Thus

[
J(u_j)
\ge
|\nabla u_j|_{L^2}^2
--------------------

C|f|*{L^2}|\nabla u_j|*{L^2}.
]

Let

[
t_j=|\nabla u_j|_{L^2}.
]

Then

[
J(u_j)\ge t_j^2-Ct_j.
]

Since (J(u_j)\to m), the values (J(u_j)) are bounded above, so (t_j^2-Ct_j) is bounded above, which implies (t_j) is bounded.

Therefore,

[
(u_j)\text{ is bounded in }H_0^1(U).
]

---

# Step 3: Extract a weakly convergent subsequence

Since (H_0^1(U)) is reflexive, there exists a subsequence (still denoted (u_j)) and some (u_0\in H_0^1(U)) such that

[
u_j\rightharpoonup u_0
\qquad\text{weakly in }H_0^1(U).
]

---

# Step 4: Use weak lower semicontinuity

The map

[
u\mapsto \int_U |\nabla u|^2,dx
]

is weakly lower semicontinuous, so

[
\int_U |\nabla u_0|^2,dx
\le
\liminf_{j\to\infty}
\int_U |\nabla u_j|^2,dx.
]

Also, the linear functional

[
u\mapsto \int_U fu,dx
]

is weakly continuous, so

[
\int_U fu_j,dx
\to
\int_U fu_0,dx.
]

Therefore,

[
J(u_0)
======

\int_U |\nabla u_0|^2,dx-\int_U fu_0,dx
]

satisfies

[
J(u_0)
\le
\liminf_{j\to\infty}
J(u_j).
]

Since (J(u_j)\to m),

[
J(u_0)\le m.
]

But (m) is the infimum, so

[
J(u_0)=m.
]

Thus the infimum is attained.

---

# Final conclusion

There exists (u_0\in H_0^1(U)) such that

[
\boxed{
J(u_0)=
\inf_{u\in H_0^1(U)}
\left(
\int_U |\nabla u|^2,dx-\int_U fu,dx
\right).
}
]

So the minimization problem has a minimizer.

---

# Why this works

The direct method always uses:

---

### 1. Minimize:

Take a minimizing sequence.

---

### 2. Compactness:

Show boundedness, then extract weak convergence.

---

### 3. Lower semicontinuity:

Pass to the limit.

---

This gives existence of a minimizer.

---

# PDE interpretation

The minimizer (u_0) is the weak solution of

[
-\Delta u=\frac{f}{2}
]

(up to constants in the functional).

Thus the variational problem corresponds to an elliptic PDE.

This variational viewpoint is central in Partial Differential Equations.

---

If you'd like, next I can show how to derive the **Euler–Lagrange equation** of this minimization problem and recover the PDE satisfied by the minimizer.

